\section{Discussion}
\label{sec:discussion}

\paragraph{Why V2 stays useful near saturation.} A benchmark that frontier systems nearly solve can look finished. We see four reasons why V2 is not. First, saturation at the frontier is itself a result: on a correctly graded split, the strongest systems fail fewer than 1\% of tasks, so small differences among them should not be read as differences in capability, on V2 or on the original split. Second, most systems that practitioners deploy are not at the frontier. Open-weight, smaller and cheaper systems still span \inklingLbPct--\kimiLbPct\% on the full set and \hardHaikuPct--\hardSonnetPct\% on Hard (Table~\ref{tab:main}), and because the residual defects we can detect affect under 1\% of tasks, a difference of a few points between such systems now reflects the systems themselves. Third, a clean set makes quantities other than the pass rate comparable. Cost, time to solution and the choice of scaffold can be compared only when every task is solvable and correctly graded: \geminiTimeouts{} of Gemini's runs reached the time limit, for instance, and mean API spend ranged from \costAstra{} to \costGemini{} per task. Fourth, a locked and fully solvable suite works as a regression test for agent products, since a failure on a task that every recent system solves is a signal worth investigating. The Hard split and the never-public set complement the full set by extending its range at the top and by guarding against contamination of the public tasks.

\paragraph{Lessons for benchmark maintainers.} Running the audit, repair and evaluation at scale taught us five lessons. (1)~\emph{Use strong agents as auditors.} The failure audit flagged \nFailureFlagsNew{} tasks that the artifact audit missed, because some defects appear only when a capable agent writes a reasonable solution that the tests reject. (2)~\emph{Repair the instruction, not the test, and validate the repair blind.} Weakening a test to fit an instruction silently changes what the benchmark measures, whereas a blind re-implementation directly checks that the edited instruction suffices; engineer review cost a median of \medianHoursPerTask{} hours per task. (3)~\emph{Never grade where the agent ran.} Agents that cannot fetch a dependency will forge one, and only a clean container rejects the forgery. (4)~\emph{Probe the lock before every run.} On the never-public set, three consecutive task exports silently dropped the no-network field and re-opened egress under unchanged harness flags (Table~\ref{tab:privdefects}); we now probe the lock inside the sandbox, under the job's exact flags, before every run. (5)~\emph{Version the benchmark and release every grade.} Defects keep surfacing as agents improve, so a benchmark needs maintenance after release, and publishing per-task grades in both grading modes lets others audit our decisions and re-score past runs.
